\documentclass[11pt]{article}

\usepackage[margin=0.7in]{geometry}
\usepackage[T1]{fontenc}
\usepackage{lmodern}
\usepackage{microtype}
\usepackage[hidelinks]{hyperref}

\pagestyle{plain}
\setlength{\parindent}{0pt}
\setlength{\parskip}{5pt}

\newcommand{\topic}[1]{\textbf{#1}\enspace}

\begin{document}

\begin{center}
    {\large EECS 221: Machine Learning for Science\par}
    \vspace{5pt}
    {\LARGE\bfseries
    Data--Compute Trade-offs\\[3pt]
    in Neural PDE Forecasting\par}
    \vspace{7pt}
    {\Large Suha Kagan Kose\par}
    \vspace{6pt}
    {\large Project Abstract\par}
\end{center}

\topic{Problem.}
Neural PDE forecasting models must be useful under limited data and
compute. We compare how accurately and efficiently three architecture
families learn fluid dynamics when training data and GPU time are limited.

\topic{Data and task.}
We use \texttt{rayleigh\_benard} from The Well~\cite{ohana2024well}: 1,750 Rayleigh--B\'enard convection trajectories across 35 (Rayleigh, Prandtl) combinations, each with 200 frames on a $512\times128$ grid, periodic horizontally and Chebyshev-spaced vertically (342\,GB). Given four previous frames, models predict the next frame's buoyancy, pressure, and horizontal and vertical velocities, then forecast further autoregressively. We keep the provided trajectory-level 80/10/10 splits and normalize every field with The Well's training-split statistics.

\topic{Models and baselines.}
We train three model families from scratch with The Well's implementations: convolutional CNextU-Net~\cite{liu2022convnext} (18.6M parameters), spectral Tensorized Fourier Neural Operator (TFNO)~\cite{kossaifi2024tfno} (19.3M), and attention-based Axial Vision Transformer (AViT)~\cite{mccabe2024mpp}, widened (hidden width 384$\rightarrow$432) to 18.6M. To check our pipeline, we first evaluate The Well's released CNextU-Net and TFNO checkpoints and compare with their published VRMSE.

\topic{Core study (${\sim}$100 GPU-hours).}
We cross two nested training-trajectory fractions (10\% and 100\%) with two seeds: $3\times2\times2=12$ runs. Each run trains for 4 hours on one A100, and we evaluate checkpoints at 1 and 4 hours, the two compute budgets. Models share subsets, targets, loss, batch size, learning-rate schedule, and evaluation horizons; an equal-budget learning-rate search on validation data precedes training.

\topic{Extensions, if time and compute allow.}
In priority order: input-noise robustness (evaluation only), a third seed, a 30\% data fraction, and a fourth, recurrent model family (ConvLSTM).

\topic{Metrics.}
Field-wise VRMSE (RMSE divided by the target field's standard deviation, as in The Well) for one-step predictions and rollouts, and steps until a predefined error threshold, with variation across seeds. Horizontal kinetic-energy spectra assess physical fidelity. We also report GPU-hours, inference latency, and peak GPU memory.

\topic{GPU estimate.}
Considering that each run uses a single GPU (e.g. NVIDIA A100  with 80\,GB VRAM), the core study requires $3\times2\times2\times4=48$ GPU-hours of training plus about 20 for profiling, tuning, baseline checks, and evaluation, with about 30 held as contingency: approximately 100 GPU-hours.

\topic{Backup plan.}
If storage, queue time, or pilot measurements show the plan is too costly, we will reduce the data size (downsampled fields or fewer trajectories) for all models equally.

\topic{Milestones for the Nov 12 check-in.}
Data pipeline validated against The Well's released checkpoints, budgets fixed from pilot timings, and first 10\%-data results with a revised plan.

\topic{Expected outcome.}
Reproducible code, accuracy--cost and data-scaling curves, and physical-fidelity comparisons showing when each architecture is most effective, including negative results.

\newpage

\begin{thebibliography}{9}
\setlength{\itemsep}{8pt}

\bibitem{ohana2024well}
Ohana, R., et al.
``The Well: a Large-Scale Collection of Diverse Physics
Simulations for Machine Learning.''
NeurIPS, Datasets and Benchmarks Track, 2024.\\
\url{https://arxiv.org/abs/2412.00568}

\bibitem{liu2022convnext}
Liu, Z., Mao, H., Wu, C.-Y., Feichtenhofer, C., Darrell, T.,
and Xie, S.
``A ConvNet for the 2020s.''
CVPR, 2022.\\
\url{https://arxiv.org/abs/2201.03545}

\bibitem{kossaifi2024tfno}
Kossaifi, J., Kovachki, N., Azizzadenesheli, K.,
and Anandkumar, A.
``Multi-Grid Tensorized Fourier Neural Operator for
High-Resolution PDEs.''
TMLR, 2024.\\
\url{https://arxiv.org/abs/2310.00120}

\bibitem{mccabe2024mpp}
McCabe, M., et al.
``Multiple Physics Pretraining for Spatiotemporal
Surrogate Models.''
NeurIPS, 2024.\\
\url{https://arxiv.org/abs/2310.02994}

\end{thebibliography}

\end{document}